% ---
% id: receiver/rem:blockmeaning
% kind: rem
% title: "Scope of the photon-number block structure"
% summary: >-
%   The block structure is in the total input photon number, which equals the
%   channel photon number only without a probe; an occupation-diagonal probe
%   state suffices to recover the channel grading. The emitted low
%   blocks decompose the POVM and are not compressions or a squashing
%   argument.
% prominence: standard
% status: accepted
% uses: [thm:click-povm]
% mentions: [def:probe]
% assumes: []
% refs: []
% provenance: newly-derived
% lean:
% ---

The difference between an operator block and a compression is the third
caveat of the sector convention
(\repodoc{spec/notation.md#24-sector-decompositions-and-block-restrictions}).

\begin{remark}\label{rem:blockmeaning}
The grading in Theorem~\ref{thm:click-povm}\,(iii) is by the \emph{total} photon
number of the declared inputs $\Inp=\InpC\sqcup\InpS$. When a probe is
declared, that is not the same as the channel photon number, and the
downstream modules need the \emph{channel} grading; recovering it is what the
occupation-diagonality of the probe states (Definition~\ref{def:probe}) is
required for, and the submodule that conditions on the probe states proves
that it suffices. When no probe is
declared ($\InpS=\varnothing$) the two gradings coincide. The emitted
low-photon blocks $\blk{M}{\leqN}_S$ are operators on
$\Hs{\Chan_{\leq N}}$ and are \emph{not} the compressions
$\Pnum{\le N}M_S\Pnum{\le N}$; emitting them is a decomposition of the
physical POVM, never a squashing argument.
\end{remark}
